\chapter{Finitud proyectiva y selector armónico dodecafásico}
\label{chap:finitud-anisotropia}
\index{contratérmino}\index{regulador externo}
\index{anisotropía}\index{CMB}\index{dodecafase}

\section{Refinamiento interno del selector armónico independiente de contraterminos externos}
\label{sec:refinamiento-sin-contraterminos}

La afirmación de que HMT no requiere renormalización fundamental distingue
tres operaciones: normalización de estados o probabilidades, refinamiento
discreto con cambio de escala y renormalización perturbativa mediante
regulador, parámetros desnudos divergentes y contraterminos dependientes del
corte. Las dos primeras son compatibles con HMT; la tercera no se introduce
como primitiva.

Sean
\[
 p_{m+1,m}:X_{m+1}\longrightarrow X_m
\]
las proyecciones sobreyectivas de un sistema inverso de espacios finitos y
sean
\[
 \mathcal O_m=C(X_m)
\]
las álgebras involutivas de observables complejos. Los morfismos de
contravariancia
\[
 p_{m+1,m}^{\ast}:\mathcal O_m\longrightarrow\mathcal O_{m+1},
 \qquad
 p_{m+1,m}^{\ast}f=f\circ p_{m+1,m},
\]
son \(\ast\)-homomorfismos unitarios coherentes. Un observable cilíndrico es
una clase compatible bajo estos morfismos.

Para medidas de probabilidad, la consistencia exacta exige
\[
 (p_{m+1,m})_\ast\mu_{m+1}=\mu_m.
\]
Para estados algebraicos positivos y normalizados
\(\omega_m:\mathcal O_m\to\CC\), la condición es
\[
 \omega_{m+1}\circ p_{m+1,m}^{\ast}=\omega_m.
\]
Cuando \(\omega_m(f)=\int f\,d\mu_m\), ambas formulaciones son equivalentes.

\begin{theorem}[Compatibilidad proyectiva y defectos sumables]
\label{thm:compatibilidad-proyectiva-defectos}
Sea \(f\in\mathcal O_m\) un observable cilíndrico y supóngase que su familia
de representantes es uniformemente acotada.
\begin{enumerate}
 \item Si las medidas o los estados son proyectivamente consistentes, sus
 expectativas son compatibles en todos los refinamientos.
 \item Si
 \[
 \left\|(p_{m+1,m})_\ast\mu_{m+1}-\mu_m\right\|_{\rm TV}
 \leq\varepsilon_m,\qquad
 \sum_m\varepsilon_m<\infty,
 \]
 entonces las expectativas forman una familia de Cauchy. La misma conclusión
 vale para estados cuando el defecto se mide en la norma dual de
 \(\mathcal O_m^\ast\).
\end{enumerate}
\end{theorem}

\begin{proof}
En el caso exacto,
\[
 \int_{X_{m+1}}f\circ p_{m+1,m}\,d\mu_{m+1}
 =\int_{X_m}f\,d\mu_m,
\]
y la identidad análoga para estados es su propia condición de consistencia.
En el caso aproximado, la dualidad entre la norma uniforme y la variación
total da
\[
 \left|
 \int f\circ p_{m+1,m}\,d\mu_{m+1}
 -\int f\,d\mu_m
 \right|
 \leq \|f\|_\infty\varepsilon_m.
\]
La acotación uniforme y la sumabilidad de los defectos convierten la suma de
incrementos en una serie convergente; por tanto, la familia es de Cauchy. La
prueba en la norma dual es idéntica.
\end{proof}

Para acciones y medidas de referencia, la condición correcta es la
consistencia de las medidas de Gibbs:
\[
 (p_{m+1,m})_\ast
 \left[
 Z_{m+1}^{-1}e^{-S_{m+1}}\,d\nu_{m+1}
 \right]
 =
 Z_m^{-1}e^{-S_m}\,d\nu_m.
\]
No basta una igualdad desnuda entre acciones, pues un mapa de comparación
podría ocultar una sustracción o reescala dependiente del corte.

La conclusión publicable queda, por ello, restringida a sus hipótesis:
donde existe una familia proyectivamente consistente de estados o medidas,
las expectativas de los observables cilíndricos uniformemente acotados son
compatibles; con defectos en variación total o norma dual sumables, forman
una familia de Cauchy. No se requiere regulador externo ni contraterminos
externos, aunque pueden existir flujos efectivos internos finitos.

En el sistema proyectivo declarado se cumplen exactamente la compatibilidad
de expectativas y el criterio de Cauchy enunciados. Una teoría cuántica de
campos completa requiere además
convergencia de correladores, positividad, unitariedad, independencia
cofinal y ausencia de dependencia del corte oculta en mapas, acciones,
medidas o parámetros. Un sector que sólo converja tras introducir parámetros
dependientes del corte refuta la versión fuerte para ese sector.

\section{Selector armónico del ciclo dodecafásico}
\label{sec:selector-armonico-c12}

Sea \(\mathsf C_{12}^{\rm cyc}\) el ciclo dodecafásico local,
\[
 \phi_k=\phi_0+\frac{2\pi k}{12},
 \qquad k=0,\ldots,11,
\]
y considérese la familia de pesos
\[
 w_k=\frac1{12}
 \left[1+2\varepsilon\cos(3\phi_k-\delta)\right],
 \qquad |\varepsilon|\leq\frac12.
\]
Entonces
\[
 w_k\geq\frac{1-2|\varepsilon|}{12}\geq0,
 \qquad
 \sum_{k=0}^{11}w_k=1.
\]

\begin{theorem}[Soporte Fourier dodecafásico]
\label{thm:soporte-fourier-dodecafasico}
La transformada discreta
\[
 \widehat w_m=\sum_{k=0}^{11}w_ke^{-im\phi_k}
\]
se anula salvo para
\[
 m\equiv0,\ \pm3\pmod{12}.
\]
Absorbida la fase de origen en \(\delta\),
\[
 \widehat w_0=1,\qquad
 \widehat w_3=\varepsilon e^{-i\delta},\qquad
 \widehat w_{-3}=\varepsilon e^{i\delta}.
\]
\end{theorem}

\begin{proof}
La transformada de la parte uniforme es el carácter trivial del ciclo. Al
escribir
\[
 2\cos(3\phi_k-\delta)
 =e^{i(3\phi_k-\delta)}+e^{-i(3\phi_k-\delta)},
\]
las sumas geométricas de los caracteres de \(\mathbb Z/12\mathbb Z\) se
anulan a menos que \(m-3\) o \(m+3\) sea múltiplo de doce. Las tres
amplitudes restantes son las indicadas.
\end{proof}

Para una corona a colatitud fija, el soporte permite únicamente \(m=0\) en
\(\ell=2\), mientras que en \(\ell=3\) pueden sobrevivir
\(m=0,\pm3\). El valor \(\varepsilon=0\) recupera la plantilla uniforme
axial; \(\varepsilon\neq0\) introduce la modulación triádica. El resultado
reconcilia ambos soportes como una envolvente común, pero no fija amplitudes
ni demuestra exclusividad: el modo \(m=0\) puede persistir también para
\(\ell=3\).

\section{Predicción de anisotropía dodecafásica: parametrización previa, contraste ciego y criterio de incompatibilidad observacional}
\label{sec:interfaz-cosmologica-c12}

La promoción a una predicción de anisotropía exige congelar antes de observar
el cielo:
\begin{enumerate}
 \item el mapa \(\mathsf C_{12}^{\rm cyc}\to S^2\) y su equivarianza;
 \item el marco celeste y el eje, o conjunto finito de ejes;
 \item \(\varepsilon,\delta,\phi_0\), la colatitud y la normalización;
 \item los \(a_{\ell m}^{\rm HMT}\), su paridad y el estadístico de
 alineamiento;
 \item máscara, primeros planos, haz, ruido y covarianza;
 \item la penalización por selección entre ejes;
 \item el contraste ciego conjunto en temperatura y polarización.
\end{enumerate}
La selección retrospectiva del eje o el ajuste a \(C_\ell^{\rm CMB}\) es
\textsc{reconocimiento externo}, no predicción. El objeto
\(\mathsf C_{12}^{\rm cyc}\) no es \(U_{12}^{\rm sgn}\), una
representación de \(SU(12)\) ni el espectro angular
\(C_\ell^{\rm CMB}\).

\begin{criterion}[Anisotropía dodecafásica]
\label{crit:anisotropia-dodecafasica}
La interfaz queda refutada si el patrón sellado es incompatible con el
contraste ciego o si eje, amplitud o modulación se seleccionan después de
abrir los datos. El teorema de soporte permanece válido aun cuando falle su
realización cosmológica.
\end{criterion}
